\documentclass[12pt,a4paper]{article}

% \usepackage[left=1.5cm,right=1.5cm,top=3cm,bottom=2cm]{geometry}
% \usepackage{titling}
% \setlength{\droptitle}{-3cm}
% \linespread{0.8}

\usepackage{savetrees}
\usepackage[colorlinks=true, urlcolor=blue, linkcolor=red]{hyperref}

\predate{}
\postdate{}
\date{}

\title{Assessing the Relevance of Search Results Obtained Through TF-IDF
Compared to Vector Embeddings}
\author{***REMOVED***}

\begin{document}
\maketitle
The aim of this project is to determine whether TF-IDF or vector search are
better at finding relevant results for queries related to computing, technology
and software. I chose this project as I have always wanted to try create a
search engine and it seems like an interesting way to apply data science with my
programming knowledge. For the data set, I created my own TF-IDF index and
vector embeddings based on \href{https://nlp.stanford.edu/projects/glove/}{GloVe
embeddings}. Although pre-trained indexes were likely available, I
decided it would be best to implement everything myself to better understand the
algorithms and methods used in vector and TF-IDF search. 

I generated the indexes with my own crawler, using websites related to the
topics as seeds. The crawler first saves the content of all the sites, then it
indexes them by generating an average vector embedding and TF-IDF values from all of the
words on the page. This data could be incomplete as some sites
may have bot blockers. For the data wrangling, I made all the tokens lowercase
and for the feature engineering, I converted the initial token counts to an
actual TF-IDF index. For the analysis, I created a set of queries
including ones that are challenging for both TF-IDF and vector search. Using these
queries, I added a page dedicated to quick evaluation of each method's results which recorded the
results to be presented after on a different page. For each query, I judged the
top three results and recorded the number of them that were relevant. These
results were later used for a per-category breakdown of where each method
struggled and succeeded, as well as the overall success of it and other
metrics. 


Improvements: substituting stopwords

Begin with the data generating process: what is the provenance of the data, how
was it collected, what limitations are known to be present, and what was its
structure at the project outset? Continue with data wrangling and feature
engineering: what steps did you take to organize the data so it was suitable
for analysis? Provide details of your analysis approach, including the type of
problem, information about key outcome variables and features, the models used
the address the question, strategies for training parameters and
hyper-parameters, and the methods used to evaluate the performance of your
approach. Present intermediate as well as final results, where appropriate.
These might cover, for example, performance metrics during training, testing
and validation, and deployment of the model to data with no known ground truth.
Try to capture the main points using visualization tools like figures or
tables. Explain results factually, with reference to concrete metrics, figures
or tables. You may save interpretation of your results to the Discussion, or
interpret them as you go with a combined Results and Discussion section.
Interpret your results in the light of your original question, if you haven’t
already done so. Discuss strengths and weaknesses, successes and failures of
your project. Compare the execution to the project plan. What could have been
improved, or what would you do differently if repeating the exercise? What
challenges arose during the project, for example data deficiencies or biases
that became apparent? What light does your project shed, if any, on the broader
area? Reflect on the ethical, societal or practical implications of your
findings. What future directions does it suggest?
\end{document}
